%==============================================================================
% ERU Symposium --- Extended Abstract
%
% Formatting follows "1. ERUS Extended-Abstract_Template.docx":
%   A4 paper, two-column body, IEEE conference style, TWO PAGES INCLUDING
%   REFERENCES. The template shows no separate Abstract block (keywords follow
%   the author block directly), so none is used here.
%
% Compile:  pdflatex extended_abstract.tex   (twice, for references)
%
% >>> BEFORE SUBMITTING: fill in every field marked TODO. <<<
%==============================================================================
\documentclass[conference,a4paper,10pt]{IEEEtran}
\IEEEoverridecommandlockouts

\usepackage[a4paper,top=0.75in,bottom=1.0in,left=0.63in,right=0.63in,
            columnsep=0.25in]{geometry}
\usepackage[T1]{fontenc}
\usepackage[utf8]{inputenc}
\usepackage{cite}
\usepackage{amsmath,amssymb}
\usepackage{graphicx}
\usepackage{booktabs}
\usepackage{textcomp}
\usepackage{xcolor}

% Ragged-bottom columns read better in a 2-page extended abstract.
\raggedbottom

\begin{document}

\title{Interleaved Frontier Exploration and Camera Coverage\\
for Target Retrieval under Short-Range Perception}

%------------------------------------------------------------------------------
% AUTHOR BLOCK
% The template is built for up to six authors, laid out left-to-right and then
% down. Delete unused \IEEEauthorblockN/A pairs; with fewer than three authors
% the remaining blocks re-centre automatically.
% TODO: confirm the spelling of every name, department, institution and city.
%------------------------------------------------------------------------------
\author{
\IEEEauthorblockN{Hiruna Malavipathirana}
\IEEEauthorblockA{\textit{Computer Science and Engineering Department} \\
\textit{University of Moratuwa} \\
Matara, Sri Lanka \\
hirunamalavipathirana.333@gmail.com}
}

\maketitle

\begin{IEEEkeywords}
autonomous exploration, coverage path planning, mobile robotics, ROS~2,
sensor footprint
\end{IEEEkeywords}

%==============================================================================
\section{Introduction}
%==============================================================================
Autonomous retrieval of objects from an unknown indoor space requires a robot to
do two things that are usually treated as one: build a map of the space, and
\emph{observe} the space with the sensor that actually detects the target. On a
low-cost differential-drive platform these two tasks are performed by sensors
with radically different footprints. A 2D~LiDAR maps a full \(360^{\circ}\)
sweep out to several metres, whereas an RGB-D camera used for object detection
covers a narrow forward cone --- on our platform a \(57^{\circ}\) horizontal
field of view (\(\pm 28.5^{\circ}\)) out to a usable depth of approximately
\(1.0\,\mathrm{m}\).

Classical frontier exploration~\cite{yamauchi} terminates when no frontier
between known-free and unknown space remains, i.e.\ when the \emph{mapping}
sensor has seen everything. This criterion says nothing about whether the
\emph{detection} sensor has ever pointed at a given patch of floor. A robot can
therefore declare exploration complete while most of the mapped floor has never
entered the camera cone, and any target lying there is missed.

This work addresses that mismatch. The \textbf{aim} is to raise the fraction of
mapped floor actually inspected by the short-range detector, without an
unacceptable increase in mission time. The \textbf{objectives} are to (i)~track,
online, which map cells the camera cone has passed over; (ii)~use that quantity
to decide \emph{when} to interrupt exploration for a systematic coverage sweep;
and (iii)~evaluate the interleaved policy against an explore-then-sweep baseline
on real hardware.

%==============================================================================
\section{Literature Review}
%==============================================================================
Frontier-based exploration~\cite{yamauchi} remains the dominant paradigm for
autonomous mapping and underlies most ROS implementations. Coverage path
planning is a separate and mature literature, surveyed by
Choset~\cite{choset} and Galceran and Carreras~\cite{galceran}, in which
boustrophedon (lawnmower) decompositions guarantee that a known region is
swept by a tool of known width.

The two are normally solved in sequence: explore to build a map, then plan
coverage over the completed map. That separation is sound when the mapping and
working footprints coincide, but fails when the detector is much
shorter-ranged, because exploration's termination criterion is blind to
detector coverage. Next-best-view and active-search methods select viewpoints
for a target sensor but generally assume the map already exists.

The gap addressed here is therefore one of \emph{scheduling}: given that both
exploration and detector coverage are necessary, and the second is far slower
per unit area, when should a robot switch between them?

%==============================================================================
\section{Materials and Methods}
%==============================================================================

\subsection{Platform and Software Stack}
Experiments used a Kobuki differential-drive base with an RGB-D camera and a
360\textdegree{} 2D~LiDAR, computing on an NVIDIA Jetson Orin Nano. The stack is
ROS~2 Jazzy with RTAB-Map~\cite{rtabmap} for SLAM and
Nav2~\cite{nav2} for planning and control, using the DWB local
planner~\cite{dwa}. Target detection uses a YOLO model in a GPU container.
The robot footprint is a measured asymmetric polygon
(\(0.36\,\mathrm{m}\) forward, \(0.22\,\mathrm{m}\) rear,
\(0.215\,\mathrm{m}\) half-width), which matters because every
inflation-derived test in Nav2 is built on an inscribed circle and cannot see
the forward overhang.

\subsection{Swept-Coverage Tracking}
A \emph{swept mask} is maintained in parallel with the occupancy grid. On every
pose update, the map cells falling inside the camera cone
(\(57^{\circ}\) horizontal FOV, \(\pm 28.5^{\circ}\), \(1.0\,\mathrm{m}\)) are marked as inspected. The
coverage deficit is then the ratio of un-swept to total known-free cells.

\subsection{Interleaving Policy}
Two arms were implemented and compared:
\begin{itemize}
  \item \textbf{Exhaustion (baseline).} Explore frontiers until none remain,
        then perform one coverage sweep --- the conventional sequential order.
  \item \textbf{Fraction (proposed).} Interrupt frontier exploration for a
        boustrophedon sweep as soon as the un-swept known-free area exceeds
        \(15\%\) of known free space, then resume exploring.
\end{itemize}
Sweep rows are spaced at \(0.86\,\mathrm{m}\), a \(10\%\) overlap on the
\(0.95\,\mathrm{m}\) swath implied by the camera cone.

\subsection{Navigation Extensions}
Two Nav2 plugins were added because stock components degraded the sweep.
First, a custom straight-line global planner is selected \emph{per goal} for
sweep rows, since the default gradient planner curves off pre-validated rows
toward locally cheaper cells. Second, a cost-aware global
planner (\texttt{SmacPlanner2D}) replaces the gradient planner for transit
goals, for the reason established in Section~\ref{sec:results}.

%==============================================================================
\section{Results and Discussion}
\label{sec:results}
%==============================================================================

\subsection{Interleaving Improves Detector Coverage}
Table~\ref{tab:arms} compares the two policy arms over hardware runs of at
least 14 minutes. The interleaved policy reached a mean final camera coverage
of \(58.2\%\) against \(42.6\%\) for the sequential baseline, an improvement of
\(15.6\) percentage points (\(37\%\) relative).

\begin{table}[t]
\caption{Final camera coverage by policy arm (hardware runs \(\geq\)14\,min)}
\label{tab:arms}
\centering
\begin{tabular}{lccc}
\toprule
\textbf{Policy} & \textbf{Runs} & \textbf{Coverage (\%)} & \textbf{Mean (\%)} \\
\midrule
Exhaustion (baseline) & 3 & 37.7--46.9 & 42.6 \\
Fraction (proposed)   & 7 & 48.6--66.4 & \textbf{58.2} \\
\bottomrule
\end{tabular}
\end{table}

The mechanism is straightforward: the baseline spends its entire exploration
phase accumulating floor the camera never inspects, and a single terminal sweep
cannot recover it within a practical mission time. Interrupting early keeps the
deficit bounded.

\subsection{Actuation Faults Masquerade as Planning Faults}
A finding with wider applicability emerged during evaluation. Missions stalled
repeatedly with the navigation stack reporting healthy operation. Bench
measurement, with commanded and achieved yaw compared against an independent
gyroscope, showed that the base driver selected its pure-rotation command mode
only when linear velocity was \emph{exactly} zero; any non-zero component
routed the command into an arc mode that scales speed about the robot centre
rather than the wheels. Because the local planner almost never emits exactly
zero linear velocity while turning, nearly every mission turn was attenuated,
losing \(84\%\) to \(91\%\) of commanded yaw. At a commanded
\(0.400\,\mathrm{rad/s}\) with \(0.0105\,\mathrm{m/s}\) of linear
velocity, the achieved/commanded ratio was \(0.16\); at
\(0.800\,\mathrm{rad/s}\) it fell to \(0.09\). Testing on a turn-radius
criterion instead of an equality test restored the ratio to \(1.01\)--\(1.03\),
while a genuine-arc control case was unchanged at \(0.78\), confirming the
change is confined to the near-pivot regime.


\subsection{Residual Stalls Are a Planning Problem}
After the actuation fix, stalls persisted. Instrumenting the controller at a
deadlock showed the robot was not in contact with anything --- the nearest
LiDAR return was \(0.41\,\mathrm{m}\) --- yet the trajectory critic rejected
all 440 sampled trajectories in 60 separate control cycles. Sampling the raw
costmap under the footprint showed it occupied inscribed-inflated cells in 54
of 64 consecutive samples: the robot habitually travelled inside the inflation
band, where trajectories toward the obstacle are correctly vetoed while an
oscillation guard suppresses the reversal that would turn it away. The cause is
that the gradient global planner weights obstacle cost weakly and produces
wall-hugging paths, which the local planner then follows faithfully. Replacing
it with a cost-aware planner that routes through the centre of free space
addressed the cause rather than the symptom.

%==============================================================================
\section{Conclusion}
%==============================================================================
The aim was met: tracking detector coverage explicitly and interleaving
boustrophedon sweeps with frontier exploration raised mean camera coverage from
\(42.6\%\) to \(58.2\%\) on hardware, confirming that map completeness is an
inadequate stopping criterion when the detection footprint is much smaller than
the mapping footprint.

Three limitations qualify these results. The comparison rests on a modest
sample (3 baseline and 7 interleaved runs) in a single arena. The navigation
extensions were introduced alongside other changes, so the contribution of each
individual component is not yet isolated. Finally, the sweep trigger threshold
was fixed at \(15\%\) rather than tuned or adapted.

Future work will isolate the planner contribution through matched control runs,
adapt the trigger threshold to observed target density, and extend coverage
bookkeeping from a binary swept mask to a confidence field that accounts for
viewing distance and angle, since detection reliability is not uniform across
the camera cone.

%==============================================================================
\section*{Acknowledgment}
The author wishes to express sincere gratitude to project supervisor Dr.~Sulochana Sooriyaarachchi, Senior Lecturer in the Department of Computer Science \& Engineering, University of Moratuwa, for her continuous guidance, invaluable insights, and technical review throughout this research. The author also gratefully acknowledges the Intellisense Lab, Department of Computer Science \& Engineering, University of Moratuwa, for providing the experimental facilities, robotic hardware platforms, and funding support essential to this work.

%==============================================================================
\begin{thebibliography}{00}

\bibitem{yamauchi} B.~Yamauchi, ``A frontier-based approach for autonomous
exploration,'' in \emph{Proc. IEEE Int. Symp. Computational Intelligence in
Robotics and Automation}, 1997, pp.~146--151.

\bibitem{choset} H.~Choset, ``Coverage for robotics --- a survey of recent
results,'' \emph{Annals of Mathematics and Artificial Intelligence}, vol.~31,
pp.~113--126, 2001.

\bibitem{galceran} E.~Galceran and M.~Carreras, ``A survey on coverage path
planning for robotics,'' \emph{Robotics and Autonomous Systems}, vol.~61,
no.~12, pp.~1258--1276, 2013.

\bibitem{rtabmap} M.~Labb\'{e} and F.~Michaud, ``RTAB-Map as an open-source
lidar and visual simultaneous localization and mapping library for large-scale
and long-term online operation,'' \emph{Journal of Field Robotics}, vol.~36,
no.~2, pp.~416--446, 2019.

\bibitem{nav2} S.~Macenski, F.~Mart\'{i}n, R.~White, and J.~Gin\'{e}s Clavero,
``The Marathon 2: A navigation system,'' in \emph{Proc. IEEE/RSJ Int. Conf.
Intelligent Robots and Systems (IROS)}, 2020, pp.~2718--2725.

\bibitem{dwa} D.~Fox, W.~Burgard, and S.~Thrun, ``The dynamic window approach
to collision avoidance,'' \emph{IEEE Robotics and Automation Magazine}, vol.~4,
no.~1, pp.~23--33, 1997.

\end{thebibliography}


%==============================================================================
% CUT FOR LENGTH --- restore if the compiled PDF has spare room on page 2.
% This was a subsection of Results and Discussion, placed after
% "Residual Stalls Are a Planning Problem". It is the most transferable
% methodological point in the work, so it is the first thing worth restoring.
%
% \subsection{A Note on Proxy Metrics}
% An intermediate result is worth recording. Extending the stall watchdog so that
% in-progress recovery behaviours were not cancelled reduced the stall rate by a
% factor of five, from \(1.26\) to \(0.23\) stalls per minute --- while making
% coverage approximately \(2.5\times\) slower, because the robot remained
% committed to unreachable goals. The stall count fell largely because that time
% ceased to be \emph{counted} as a stall. Coverage-versus-time, not stall count,
% is the metric that reflects the objective.
%
% Also cut: a table of achieved/commanded yaw before and after the driver fix.
% Its numbers now appear inline in "Actuation Faults Masquerade as Planning
% Faults"; the full five-row table is in week_8.md.
%==============================================================================

\end{document}
